
\section{Discussion}
\label{sec:discussion}

\paragraph{Why per-DOF conformal features fail.}
The negative result on conformal conditioning (Section~\ref{sec:ablation}) warrants explanation.
% exp_id: perdof_eps0.1_object → 11.0%, obs_only_eps0.1_object → 19.0%
At $K{=}10$ samples in $d{=}6$ dimensions, the dimension-to-sample ratio $d/K = 0.6$ means that even with Ledoit--Wolf shrinkage~\cite{ledoit2004well}, per-DOF statistics---standard deviations, conformal radii---carry substantial estimation noise.
The frozen observation embedding $\mathbf{z}_{\mathrm{obs}} \in \mathbb{R}^{990}$ already encodes the geometric information these features attempt to capture: object identity, spatial layout, and gripper-object relationship.
Adding noisy low-dimensional statistics to a rich representation does not improve the CRM's predictions; it degrades them.

A subtler issue is the gap between offline and closed-loop metrics.
All CRM variants achieve low per-step MSE on held-out demonstration data, yet offline prediction accuracy does not predict closed-loop success.
In closed-loop execution, compounding error~\cite{ross2011dagger} shifts the state distribution away from training data, and the noisy conformal features amplify this drift---occasional erroneous correction directions, even when small in magnitude, can trigger cascading failures in individual episodes.
This challenges the intuition from controlled synthetic settings that richer uncertainty features improve corrections, and raises an open question: would larger $K$ (e.g., $K{=}50$) yield reliable enough covariance estimates for per-DOF features to help?
We cannot test this due to the $K$-fold increase in VLA forward passes per control step.

\paragraph{When does action-level correction help?}
% exp_id: obs_only_eps0.1_object → 19.0%, argmin_K10_object → 12.0%
% exp_id: obs_only_eps0.1_spatial → 35.0%, argmin_K10_spatial → 34.0%
The +7.0\,pp gain on LIBERO-Object versus the null result on LIBERO-Spatial (Table~\ref{tab:main}) reveals a structural distinction.
Object perturbations (novel object instances) create systematic action biases: the VLA produces actions with consistent directional errors that a learned module can identify and compensate from observation features alone.
Spatial perturbations (layout changes) require higher-level re-planning---the actions may be locally reasonable but directed at the wrong goal location.
Action-level residuals are effective for fine-grained motor adjustments, not for strategic re-planning.
This suggests that the benefit of post-hoc correction depends on whether the distribution shift manifests as correctable motor error or as a planning-level failure.

\paragraph{Limitations.}
We identify seven limitations, each of which qualifies the scope of our findings:
\begin{enumerate}[nosep]
  \item \emph{Simulation only.} All experiments use LIBERO~\cite{liu2023libero}; real-robot transfer is untested.
  \item \emph{Single VLA backbone.} We evaluate only SmolVLA~\cite{shukor2025smolvla} (450M parameters). Whether ACR-VLA transfers to larger VLAs (e.g., $\pi_0$, OpenVLA) is unknown.
  \item \emph{Limited benchmarks.} Two LIBERO suites (20 tasks total); other manipulation benchmarks and perturbation types are not evaluated.
  % exp_id: obs_only per-task variance — some object tasks underperform argmin
  \item \emph{Per-task variance.} On individual LIBERO-Object tasks, obs\_only occasionally underperforms the Argmin baseline. The method does not universally help.
  \item \emph{Magnitude bound requires tuning.} We use $\varepsilon{=}0.1$ throughout, but the optimal $\varepsilon$ likely varies across VLA architectures and task domains.
  % exp_id: seed42 only; seed43 pending
  \item \emph{Single-seed results.} All primary results use seed 42. Seed 43 replication is pending; we make no statistical significance claims.
  % exp_id: obs_only_eps0.1_object → 19.0%
  \item \emph{Absolute performance.} 19.0\% success on LIBERO-Object is far from deployment-ready. ACR-VLA reduces the gap relative to selection-only but does not close it.
\end{enumerate}


\section{Conclusion}
\label{sec:conclusion}

% exp_id: obs_only_eps0.1_object → 19.0%, argmin_K10_object → 12.0%
% exp_id: obs_only_eps0.1_spatial → 35.0%, argmin_K10_spatial → 34.0%
We presented ACR-VLA, a sample-then-correct pipeline that wraps a frozen flow-matching VLA with multi-sample conformal selection followed by a lightweight residual correction module (${\sim}$330K parameters, ${\sim}$3\,ms overhead).
On LIBERO-Object, observation-conditioned CRM achieves 19.0\% task success versus 12.0\% for selection alone (+7.0\,pp).
On LIBERO-Spatial, correction provides no measurable benefit over selection (best CRM within 1\,pp of the selection baseline), indicating the approach is effective for fine-grained motor correction under object-level distribution shift but not for spatial re-planning.

% exp_id: perdof_eps0.1_object → 11.0%, obs_only_eps0.1_object → 19.0%
Our central negative result is that per-DOF conformal features do not improve correction over observation-only conditioning, despite providing intuitively richer uncertainty information.
At the $K{=}10$ sample budget, covariance estimation noise outweighs the information gain.
This is a cautionary finding for the broader agenda of uncertainty-conditioned action correction: the gap between offline prediction accuracy and closed-loop task success remains a fundamental challenge.

Concrete next steps include: (i)~validation across multiple VLA backbones and scales, (ii)~real-robot transfer experiments, (iii)~adaptive magnitude scheduling (learned $\varepsilon$ conditioned on observation uncertainty), and (iv)~higher-$K$ sampling to revisit whether per-DOF conformal features become useful with better covariance estimates.
